\section{A supplementary Galerkin certificate for the fixed datum}
\label{n3:section}

This section records a computer-assisted estimate complementary to the
Picard convergence result. All norms use normalized measure on
$\mathbb T^3=(\mathbb R/2\pi\mathbb Z)^3$, the viscosity is one, and fields
are mean zero and divergence free. Write $A=-\Delta$,
$B(a,b)=\mathbb P((a\cdot\nabla)b)$, and let $P_M$ retain the nonzero
Fourier modes with $|k|\le M$. Denote the exact Galerkin solution from
$P_MW$ by $u_M$, where $W=(\sin z,\sin x,\sin y)$, and put $v=u_3$.

\paragraph{Supplementary conclusion.}
The certificate and estimates below imply a unique global smooth solution
$u$ from $W$ and
\begin{align}
 \|\nabla(u_M-v)(t)\|_2^2&<10^{-3}e^{-t/2}
       &&(M\ge3),\label{n3:error}\\
 \sup_M\sup_{t\ge0}\|\nabla u_M(t)\|_2^2&<14,
 &B_*:=\sup_M\int_0^\infty\|\nabla u_M(t)\|_2^4\,dt
       &<\frac{21}{2}.\label{n3:bounds}
\end{align}
The error bound also holds for the PDE limit $u$. The corresponding norm
bound is $\|\nabla(u-v)(t)\|_2<10^{-3/2}e^{-t/4}$.
Separately, if $h\in H^1$ is mean zero and divergence free and
$\|\nabla h\|_2\le1/2500$, the datum $W+h$ generates a unique global
strong solution, smooth for positive time, with
\begin{equation}
 \|\nabla(u-v)(t)\|_2^2<\frac1{500}e^{-t/2}.
 \label{n3:ball}
\end{equation}
Smooth perturbations retain smooth initial regularity. This ball has no
Fourier-support restriction. The fixed-$W$ energy estimates
\eqref{n3:bounds} are not assigned unchanged to perturbed initial data.

\subsection{The error estimate and scalar comparison}
\label{n3:error-proof}

The symbols $G(t)$ and $L(t)$ below are scalar functions local to this
section; the scalar tail coefficient $H=\sqrt{660E}$ is defined later.
They are distinct from the finite Fourier fields and profile notation
of the preceding sections.

For a smooth reference field $v$ with residual
$r=v'+Av+B(v,v)$ and a smooth exact solution on its interval of existence,
set $e=u-v$ and
$z=\|\nabla e\|_2^2$, $D=\|\Delta e\|_2^2$, and
$G=\|\nabla v\|_{L^\infty}$, using the Frobenius norm pointwise.
The vector/tensor Sobolev estimate needed here is
\begin{equation}
 \|f\|_6\le C_S\|\nabla f\|_2,
 \qquad C_S=25/3,
 \label{n3:sobolev}
\end{equation}
for componentwise mean-zero fields. To check its normalization, on a
period-$L$ torus coordinate variation gives, for nonnegative scalar $a$,
$a\le g_i$, where $g_i$ is the coordinate average of
$a+(L/2)|\partial_i a|$. Repeated Cauchy--Schwarz applied to
$a^{3/2}\le(g_1g_2g_3)^{1/2}$, followed by the arithmetic--geometric mean
inequality, yields
\[
 \|a\|_{3/2}\le\|a\|_1+\frac L6\sum_i\|\partial_i a\|_1.
\]
Apply this to $a=|f|^4$ and use H\"older and
$\sum_i\|\partial_i f\|_2\le\sqrt3\|\nabla f\|_2$ to obtain
$\|f\|_6\le\|f\|_2+(2L/\sqrt3)\|\nabla f\|_2$.
At $L=2\pi$, Poincar\'e and
$1+4\pi/\sqrt3<25/3$ prove \eqref{n3:sobolev}.
This argument applies directly to the Euclidean magnitude of vector or
tensor fields. In particular it applies to $\nabla e$, with
$\|D^2e\|_2=\|\Delta e\|_2$.

Testing the error equation with $-\Delta e$, integrating the transport
term once by parts, and using mean-zero Poincar\'e gives, initially with
viscosity $\nu>0$,
\[
 \frac12z'+\nu D
 \le Gz+G\sqrt{zD}
       +C_S^{3/2}z^{3/4}D^{3/4}+\|r\|_2\sqrt D.
\]
Allocate $\nu D/4$ to each of the last three terms. The cubic Young
estimate has remainder $27C_S^6z^3/(4\nu^3)$. Therefore
\begin{equation}
 z'+\frac\nu2D
 \le(2G+2G^2/\nu)z+\frac{27C_S^6}{2\nu^3}z^3
       +\frac2\nu\|r\|_2^2.
 \label{n3:h1-inequality}
\end{equation}
Set $\nu=1$, $f=2\|r\|_2^2$, and define
\begin{align}
 \Lambda(t)&=\int_0^t(2G+2G^2)\,ds,&
 L(t)&=\Lambda(t)-t/2,\notag\\
 \eta&=d^2+\int_0^\infty e^{-L}f\,dt,&
 I&=\int_0^\infty e^{2L}\,dt,\qquad
 \kappa=27C_S^6\eta^2 I,\label{n3:gate-def}
\end{align}
where $d^2=z(0)$. Since $D\ge z$, $w=e^{-L}z$ satisfies
$w'\le(27C_S^6/2)e^{2L}w^3+e^{-L}f$.
The remaining-forcing function
$V(t)=w(t)+\int_t^\infty e^{-L(s)}f(s)\,ds$ obeys
$V(0)=\eta$, $V\ge w\ge0$, and
$V'\le(27C_S^6/2)e^{2L}V^3$.
Bernoulli comparison, whenever $\kappa<1$, gives
\begin{equation}
 z(t)\le
 \frac{e^{L(t)}\eta}
 {\sqrt{1-27C_S^6\eta^2\int_0^t e^{2L(s)}\,ds}}
 \le\frac{e^{\Lambda(\infty)}\eta}{\sqrt{1-\kappa}}e^{-t/2}.
 \label{n3:scalar-bound}
\end{equation}

\subsection{Exact algebra, interval validation, and the infinite tail}
\label{n3:certificate}

The finite reference solves $v'+Av+P_3B(v,v)=0$ from $W$, so
$r=(1-P_3)B(v,v)$. The supplied exact rational data give a ten-dimensional
real invariant basis. Its Laplacian eigenvalues and squared norms are
\[
 (1,2,5,5,6,6,8,9,9,9),\qquad
 (3/2,3/4,3/16,15,1/8,12,6,1566,87/8,3/2).
\]
The first basis vectors are $W$ and
$F=(\sin y\cos z,\sin z\cos x,\sin x\cos y)$, using the notation of Section~\ref{main:reference}.
Exact Fourier convolution of all 55 unordered basis pairs verifies full
closure within the spherical cutoff, 133 nonzero quadratic ODE
coefficients, and the entire exterior residual. There are 52 nonzero
exterior pair fields, with 648 signed exterior modes, whose squared norm
has 285 quartic terms. The retained support $\mathcal K$ satisfies
\begin{equation}
 |\mathcal K|=108,
 \qquad\sum_{k\in\mathcal K}|k|^2=660.
 \label{n3:support}
\end{equation}
The same exact audit verifies reality, divergence freedom, basis
orthogonality, the initial datum, and cubic cancellation in the energy law.
The basis, coefficients, and residual fields are part of the accompanying
machine-readable proof data; the full coefficient lists are not printed here.

The accepted enclosure uses integer intervals with denominator $2^{224}$,
step $h=1/128$, and Taylor order $p=16$ for 1024 steps through $S=8$.
At each step strict inclusion
$X+[0,h]F(\mathcal B)\subset\operatorname{int}\mathcal B$
provides a continuous Picard enclosure. Taylor coefficients at $X$ give
the initial derivatives, and the same recurrence on $\mathcal B$ gives
universal derivative remainders. For the dissipative diagonal component
$x_k'=-\lambda_kx_k+N_k(x)$, endpoint propagation uses
\[
 e^{-\lambda_kh}X_k+\sum_{n=0}^{p-1}w_{\lambda_kn}N_{k,n}(X)
       +w_{\lambda_kp}N_{k,p}(\mathcal B),\qquad
 w_{\lambda n}=\int_0^h e^{-\lambda(h-s)}s^n\,ds.
\]
Alternating rational series enclose the exponentials and weights; integer
square roots round outwards. Ordinary Taylor bounds enclose the whole
time slab. Intersections of proved enclosures are used, never midpoint
replacement; failed inclusion or an empty intersection rejects the run.

On these full slabs, $\|r\|_2^2$ is enclosed both by the exact quartic
polynomial and by a midpoint value with an $L^2$ radius bound. For $G$, the
first two basis fields satisfy
$\|\nabla(aW+bF)\|_\infty=\sqrt{3(a^2+b^2)}$; the remaining fields
use their exact Fourier triangle bounds. For example, the first row of
the first identity is bounded by
$b^2\cos^2y\cos^2z+(a\cos z-b\sin y\sin z)^2\le a^2+b^2$;
cyclic summation and equality at the origin prove the formula.
Full-slab rectangle bounds integrate $\Lambda$, $\eta$, and $I$.

Let $E$ enclose $\|v(S)\|_2^2$. The exact Galerkin energy law implies
$\|v(S+\tau)\|_2^2\le Ee^{-2\tau}$. By \eqref{n3:support},
\[
 G(S+\tau)\le He^{-\tau},\quad H=\sqrt{660E},\qquad
 \|r(S+\tau)\|_2^2\le972E^2e^{-4\tau}.
\]
Here the latter bound follows from
$\|v\|_\infty\le\sqrt{108}\|v\|_2$,
$\|\nabla v\|_2\le3\|v\|_2$, and contraction of the residual projection.
Thus, writing $\Delta=2H+H^2$,
\begin{equation}
 \Lambda(\infty)\le\Lambda(S)+\Delta,\qquad
 \eta_{\rm tail}\le\frac47\,972E^2e^{-L(S)},\qquad
 I_{\rm tail}\le e^{2L(S)+2\Delta}.
 \label{n3:tails}
\end{equation}
The finite Galerkin energy law ensures existence throughout this tail;
no unknown global PDE solution is assumed here.

For $d=0$, the reproduced finite certificate and these tails give
\begin{equation}
 \eta<\frac{587}{10^9},\qquad I<156400,\qquad
 \Lambda(\infty)<\frac{176}{25},\qquad
 \kappa\le0.486226878267<\frac{49}{100}.
 \label{n3:numeric-bounds}
\end{equation}
The displayed decimal is an outward upper rounding. Exact rational
arithmetic also verifies
$27(25/3)^6(587/10^9)^2\,156400=134726479/276480000<49/100$.
A positive rational exponential series with a geometric remainder proves
$e^{176/25}<1150$. Consequently
$1150(587/10^9)/\sqrt{51/100}<1/1000$.
The release replay documented in Appendix~\ref{cert:section} checked all recorded interval quantities against the
archived result exactly, except for wall-clock duration; the exact algebra
and rational-margin outputs also agreed. This is a reproduction using the
same interval implementation, not an independent implementation or a
proof-assistant formalization.

\subsection{All cutoffs, the PDE limit, and the perturbation ball}
\label{n3:pde}

For each $M\ge3$, both exact finite Galerkin flows exist globally by
their energy identities. Since $v$ belongs to the $M$-mode space,
$e_M=u_M-v$ obeys \eqref{n3:h1-inequality} with residual $P_Mr$ and
$d=0$. Orthogonal projection does not increase its $L^2$ norm, so the same
certificate applies uniformly in $M$. Equations
\eqref{n3:scalar-bound}--\eqref{n3:numeric-bounds} prove
\eqref{n3:error}. Moreover,
\[
 \|v(t)\|_2^2\le\frac32e^{-2t},\qquad
 \|\nabla v(t)\|_2^2\le\frac{27}{2}e^{-2t}.
\]
Writing $X_M=\|\nabla u_M\|_2^2$ gives
\[
 \sqrt{X_M(t)}<\sqrt{27/2}\,e^{-t}+\sqrt{1/1000}\,e^{-t/4},
 \qquad \alpha:=(\sqrt{27/2}+\sqrt{1/1000})^2<14.
\]
For $M<3$, the direct spectral energy estimate gives
$X_M(t)\le(27/2)e^{-2t}$ instead. The exact energy identity gives
$\int_0^\infty X_M\,dt\le\|P_MW\|_2^2/2\le3/4$ for every cutoff.
Thus $\sup_{M,t}X_M\le\alpha<14$ and
$B_*\le3\alpha/4<21/2$, including a strict uniform margin.

Retaining the positive $D/2$ term in \eqref{n3:h1-inequality} and
integrating yields a uniform $L^2(0,T;H^2)$ bound for $e_M$, since the
error estimate, the finite certificate, and \eqref{n3:tails} control all
terms on the right. Together with the reference bounds this gives uniform
$L^\infty(0,T;H^1)\cap L^2(0,T;H^2)$ control of $u_M$.
The equation bounds $\partial_tu_M$ in $L^2(0,T;L^2)$; explicitly,
\[
 \|(u_M\cdot\nabla)u_M\|_2^2
 \le C_S^3 X_M^{3/2}\|\Delta u_M\|_2
 \le\frac{C_S^3}{2}(X_M^3+\|\Delta u_M\|_2^2).
\]
Compactness as in \cite[Corollary~4]{Simon} and strong uniqueness as in
\cite[Proposition~2.2]{Tao2009} therefore produce a
global solution in $C([0,T];H^1)\cap L^2(0,T;H^2)$ for every finite $T$.
Lower semicontinuity and time continuity transfer the reference-error
bound; smooth initial data and parabolic regularity give global smoothness.
This passage uses standard local existence, continuation, and compactness
theory, in addition to the finite certificate.

For the perturbation ball take $d^2\le1/2500^2$, or its projected value
for $u_M(0)=P_M(W+h)$. The same reference and residual give
\[
 \eta_h<\frac{747}{10^9},\qquad
 \kappa_h<27(25/3)^6(747/10^9)^2\,156400
   =\frac{8080797}{10240000}<\frac45.
\]
Since $[1150(747/10^9)]^2/(1-4/5)<(1/500)^2$,
\eqref{n3:scalar-bound} proves \eqref{n3:ball}; the same uniform
Galerkin construction applies. Its initial energy differs from that of
$W$, so this conclusion is stated separately from \eqref{n3:bounds}.

\paragraph{Reproducibility and scope.}
The accompanying files in \path{checks/N3/} of \path{Code.zip}, namely
\path{n3_exact_algebra_intermediate.json},
\path{n3_algebra_audit.py}, \path{n3_interval_core.py},
\path{n3_validated_run.py}, and \path{n3_check_margins.py}
contain the complete finite data and executable checks.
The runner \path{run_replay.py} and the release record
\path{replay_receipt.json} in the same directory identify the
reproduction and exact source hashes. Correctness still depends on
these programs, their exact-arithmetic execution, and the written analysis
above. The result concerns the fixed datum and the stated neighborhood;
it implies neither arbitrary-data three-dimensional regularity nor a
new precedence claim.
